\guidechapter{Why more calls can still mean less work}{Count the size of each job as well as the jobs.}

Suppose a width-doubling step makes three smaller calls, each on half the
data volume. Three calls sounds worse than two, but it does not mean three
full-size jobs. Their total visited volume is $3\cdot(V/2)=1.5V$.
At the next depth there are nine calls on volume $V/4$ each, totaling
$2.25V$. The circle areas in the tree encode those shrinking volumes.

\explfigure{19-recursion-ledger}{The growth factor is calls times relative volume.}{
This toy tree uses $m=2$, $W=2$, $s=3$. Its first three levels visit
$V$, $1.5V$, and $2.25V$. The right chart compares this with doubling at
every level. Calls can run sequentially and reuse storage; the total visited
volume describes work, not peak memory.}

After $j$ levels, there are $3^j$ calls and each has volume $V/2^j$.
Multiplying gives $V(3/2)^j$. Meanwhile the width factor is $2^j$.
To express the work as a power of width, solve $2^q=3/2$. This gives
$q=\log_2(3/2)$, about $0.585$, below one. The corresponding toy cost
grows more slowly than the width itself.

\paragraph{What the real proof must add.}
The tree charges only the recursive contributions shown. A real algorithm
also scans, copies, arranges, and restores records at its nodes. It must bound
that overhead, handle widths that do not fit the tree exactly, and keep one
fixed machine description. The source reserves strict slack in its recurrence
inequality for these tasks. Its numerical saving is much smaller than the
one drawn here.

\paragraph{A diagnostic question.}
What if every one of the three children actually processed the whole volume
$V$? The multiplier would be three, not $3/2$, and this argument would no
longer beat a factor of two. This is why the paper's division by $W$ is an
implementation obligation rather than a cosmetic normalization.

\takeaway{A recurrence counts weighted work. The weight is the data volume
of each child, including the effects of the required layout and workspace.}
